\section{个人知识库与自然语言偏好管理}

\subsection{多类型资料解析与安全边界}

自由行用户的行前准备散落在长文游记、景区官方 PDF 指南、网页攻略与收藏夹中。知识库服务（\texttt{knowledge/}）坚持“\textbf{本机提取、串行可控、失败明确、彻底可删}”四项原则，支持四类主动导入方式：纯文本粘贴、文件上传（TXT / Markdown / 可提取文字的 PDF）、公开网页 URL 以及收藏链接。

在解析与安全边界上，系统建立了完整的防护网：
\begin{itemize}
  \item \textbf{明确容量与资源上限（R19）}：单实例默认上限 100 份资料、单份文件上限 20 MiB、PDF 上限 200 页、单份提取正文上限 200 万字符；超限时在接口层（\texttt{409}/\texttt{413}）或解析层明确拒绝，失败资料标记为 \texttt{failed} 且绝不向索引写入半成品片段。
  \item \textbf{扫描件识别提示（R07）}：对纯图片无文字层的扫描 PDF，解析器明确判定为 \texttt{pdf\_no\_text} 并返回可读引导（“PDF 没有可提取文字，请提供文字版本或粘贴正文”），零片段入库，避免静默入库空文档误导检索。
  \item \textbf{公网网页抓取 SSRF 防护}：导入网页 URL 或收藏链接时，\texttt{public\_web.py} 校验协议、禁用凭据 URL、解析目标主机地址并拒绝回环地址、私有网段与链路本地地址，且逐跳验证每一次 HTTP 重定向，防止恶意或误填链接探测本机服务。
  \item \textbf{独立解析进程与中断恢复}：解析以受控子进程执行并施加超时与内存边界；若服务在导入长文或大 PDF 时被强制终止，重启后未完任务自动标记为 \texttt{failed/import\_interrupted}，无残留脏索引，用户可直接在界面一键“重新导入”。
\end{itemize}

\subsection{本机中文检索、按需关联与彻底删除}

完成文字提取后，系统按约 800 中文字符、相邻 100 字重叠的滑动窗口切分文本，保留原始 PDF 页码或文本段落号定位，并使用 \texttt{jieba} 中文分词同时处理入库正文与查询词，写入 SQLite FTS5 全文表。相比引入外部向量数据库与嵌入模型，这一方案实现了毫秒级本地响应（100 份容量压测下中位数 5.29 ms），且对“西安”“碑林”等高频中文短词召回稳定、可精确定位原文出处。

为了兼顾规划相关性与隐私最小化（R08），系统在 Schema v6 中实现了\textbf{行程级参考资料关联}：
\begin{enumerate}
  \item 用户可在新建或编辑行程的表单弹窗中直接勾选已有资料，或现场粘贴文字、上传文件、导入链接并即时关联到当前行程版本；取消关联仅解除本次行程引用，原资料仍保留在资料库中供后续选用。
  \item 规划期间仅检索当前行程显式关联且目的地相关的资料片段，按 6,400 字符上下文预算裁剪后仅发送给规划模型，绝不送往高德地图或联网搜索。
  \item 每次模型决策、格式修复与独立审核前都会通过 \texttt{live\_chunks} 重新核查关联片段是否仍存在；用户一旦在资料库删除某份资料，系统在事务中同步清除片段记录与 FTS5 倒排项并移除本地原始文件目录，确保后续检索与规划绝不再读取已删内容。
\end{enumerate}

\subsection{自然语言旅行偏好与优先级闭环}

早期的结构化键值偏好表单难以自然表达自由行用户的真实习惯（如“上午九点半后再出门”“逛博物馆想留足细看碑刻的时间”“午晚餐不赶网红店”）。在 Schema v6 升级中，系统将偏好档案简化为一段\textbf{自然语言原文（\texttt{notes}）}，旧键值数据自动平滑迁移为逐行文本，在前端提供简洁的查看、编辑保存与二次确认清空入口，并以 \texttt{expected\_revision} 乐观锁拦截并发旧版本覆盖（返回 \texttt{409 revision\_conflict}）。

在规划决策链条中，偏好严格服从三项规则约束（R01--R04）：
\begin{itemize}
  \item \textbf{本次优先于长期}：当用户在某次行程中明确提出与长期习惯冲突的诉求（例如长期“不逛夜景商业街”，本次“必须安排大唐不夜城夜游”），规划引擎优先满足本次明确要求，同时保持数据库中的长期偏好原文与版本号纹丝不动。
  \item \textbf{辅助而非硬约束}：长期偏好只在不突破抵离时间、场馆闭馆与全团预算等硬约束的前提下引导节奏与选点，缺失普通偏好时系统采用可见假设直接规划，不反过来弹窗打扰用户。
  \item \textbf{无多余匹配理由}：规划与介绍生成提示词均明确禁止输出“根据您的长期偏好为您安排……”式文本，让行程卡片回归景点本身的文化与体验价值。
\end{itemize}
